\section{Experimental Protocol}
\label{sec:protocol}

\begin{figure}[t]
\centering
\begin{tikzpicture}[
  node distance=0.55cm and 0.4cm,
  role/.style={draw, rounded corners=3pt, minimum width=2.6cm, minimum height=0.6cm,
               align=center, font=\small},
  train/.style={role, fill=green!15},
  cal/.style={role, fill=orange!15},
  evalneg/.style={role, fill=blue!12},
  evalpos/.style={role, fill=red!12},
  locked/.style={role, fill=gray!20, dashed},
  arrow/.style={-{Stealth[length=4pt]}, thick},
  label/.style={font=\scriptsize}
]

\node[train]   (tr)  {Train-Normal\\$\mathcal{E}_{\text{train}}$};
\node[cal, right=1.2cm of tr]  (cal) {Calibration-Normal\\$\mathcal{E}_{\text{cal}}$};
\node[evalneg, right=1.2cm of cal] (evn) {Eval Normal-Neg.\\$\mathcal{E}_{\text{eval}}^{-}$};
\node[evalpos, right=1.2cm of evn] (evp) {Eval Buggy-Pos.\\$\mathcal{E}_{\text{eval}}^{+}$};
\node[locked, below=0.7cm of evn, xshift=1.5cm] (lock)
  {Locked Test\\(not materialized)};

% Separation constraint
\draw[arrow, dashed, gray] (tr) -- node[above, label]{source $\cap$ pair\\$= \emptyset$} (cal);
\draw[arrow, dashed, gray] (cal) -- node[above, label]{source $\cap$ pair\\$= \emptyset$} (evn);
\draw[arrow, dashed, gray] (evn) -- node[above, label]{} (evp);

% Usage arrows
\draw[arrow] (tr.south) -- ++(0,-0.4) -| node[below, near start, label]{Fit (if any)} (tr.south);
\draw[arrow] (cal.south) -- ++(0,-0.4)
  node[below, label]{Set $\tau$ (Eq.~\ref{eq:threshold})};
\draw[arrow] (evn.south) -- ++(0,-0.4)
  node[below, label]{Compute metrics};
\draw[arrow] (evp.south) -- ++(0,-0.4)
  node[below, label]{Compute metrics};
\draw[->, dashed, gray!60, thick] (lock.west) -- ++(-0.5,0)
  node[left, label, gray]{Not scored};

\end{tikzpicture}
\caption{Leakage-aware split protocol. Source, pair, and episode identifiers are
assigned to roles \emph{before} window generation; no identifier crosses a role boundary.
Train-normal episodes are the only input for any learned component fit.
Calibration-normal episodes set the detection threshold on normal scores only.
Locked test data is neither materialized nor scored; all reported metrics are
non-locked validation evidence.}
\label{fig:protocol}
\end{figure}

\subsection{Role Separation and Leakage Controls}

Sources, suspected pairs, and episodes are assigned to roles before window generation. Train-normal
episodes are the only allowed input for train-dependent fitting. Calibration uses designated normal
episodes, and evaluation contains held normal-negative and buggy-positive episodes. The main
TempGlitch follow-up has zero cross-role overlap by \texttt{source}, \texttt{pair\_id}, and
\texttt{source\_episode\_id}. This grouping matters because a normal/buggy pair or derivative
capture split across roles could make evaluation easier without demonstrating transferable
anomaly detection~\cite{kaufman2011leakage}. Figure~\ref{fig:protocol} illustrates the
role structure.

\subsection{Main TempGlitch Follow-up}

The frozen pair-disjoint, non-locked TempGlitch follow-up contains 30,549 windows. Two normal
episodes define the threshold:
\texttt{Godot\_Blinking\_Normal\_106} and
\texttt{Godot\_Frozen\_Animation\_Platformer\_Normal\_107}. The evaluation support is exactly 34
episodes: 12 normal-negative and 22 buggy-positive. All 60 recorded comparison rows use this same
support tuple. TempGlitch supplies binary video labels; no temporal span metric is evaluated.

\subsection{Secondary R5-XGame Evidence}

R5-XGame is reported separately on a frozen non-locked split with 12 calibration-normal, 12
normal-negative, and 60 buggy-positive episodes. Within that bundle, all compared methods share
the same support. Its distribution, class balance, calibration support, action mode, and training
lineage differ from TempGlitch, so absolute metrics are not interpreted as a cross-benchmark
comparison.

\subsection{Metrics and Uncertainty}

We report AUROC, AUPRC, F1, precision, recall, balanced accuracy, and FPR@95TPR. TempGlitch
confidence intervals use 1,000 pair-grouped bootstrap replicates for AUROC and F1. R5-XGame uses
the corresponding validated group-bootstrap intervals from its evidence bundle. The operating
threshold is fixed from normal calibration scores before evaluation.

Validators enforce
\texttt{validation\_buggy\_used\_for\_fit\_select=false},
\texttt{locked\_test\_materialized=false}, and
\texttt{locked\_test\_scored=false}.
